\section{Introduction}
\label{sec:intro}

Vision-Language Models (VLMs) have achieved remarkable performance across a wide
range of perception and reasoning tasks~\cite{liu2023visual,bai2025qwen,team2023gemini}.
Yet they remain prone to systematic failures that are difficult to diagnose or correct
post-hoc. When a VLM answers a question about an image incorrectly, the failure may
stem from two distinct sources: (1)~a \emph{reasoning error}---incorrect inference
from accurate perceptual inputs; or (2)~a \emph{premise error}---reasoning from false
beliefs about what the image contains. Standard evaluation does not distinguish these
two failure modes, and standard training procedures do not explicitly enforce that
intermediate beliefs be grounded in visual evidence.

This paper asks a deceptively simple question: \emph{what fraction of VLM errors
are caused by false visual premises, and can correcting them at inference time recover
the correct answer?}

\paragraph{The Epistemic Failure Hypothesis.}
We hypothesize that a substantial portion of VLM failures are \emph{epistemic} in
nature---the model's answer is wrong because it begins reasoning from an incorrect
perception of the scene, not because its reasoning chain is flawed.
If this holds, then selectively correcting false beliefs at inference time, without any
weight updates, should flip wrong answers to correct.
This is a falsifiable, causal claim: under a random-ablation control (removing
beliefs without targeted correction), we would expect no such systematic improvement.

\paragraph{\bci{}: Belief-Constrained Inference.}
We propose a structured inference-time loop operating in five stages:
\begin{enumerate}[leftmargin=1.5em,itemsep=1pt,topsep=2pt]
  \item \textbf{Belief Externalization.} Before producing a final answer, the VLM
    emits 5--10 atomic, falsifiable visual claims about the image.
  \item \textbf{Claim Extraction.} Claims are parsed into typed atomic predicates
    (object, attribute, relation, action, count).
  \item \textbf{Claim Verification.} Each claim is checked against structured visual
    evidence (GQA scene graphs), yielding
    \textsc{Supported}, \textsc{Contradicted}, or \textsc{Uncertain}.
  \item \textbf{Belief Revision.} Contradicted claims are replaced with
    ground-truth facts (E1), removed (E3), or selectively corrected via a
    confidence gate (E12b).
  \item \textbf{Constrained Re-reasoning.} The model re-answers using only the
    revised belief set.
\end{enumerate}

\paragraph{Key Findings.}
On GQA under Qwen2.5-VL-7B~\cite{bai2025qwen}, 33.3\% of all incorrect answers
involve at least one contradicted belief, rising to 50\% for logical/relational
questions. Full premise replacement (E1) yields $+18.0$ pp at $n{=}500$
($p{=}1.5{\times}10^{-14}$) and $+16.2$ pp at $n{=}2000$
($p{=}1.78{\times}10^{-48}$). Across both scales, flip-to-correct rates exceed
21\%, while flip-to-wrong rates remain below 6\%.
Targeted correction outperforms random belief deletion by $+35$ pp, confirming
that the effect is claim-specific and causal.
A practical, compute-bounded gated policy (E12b) achieves $+3.45$ pp at $n{=}2000$
($p{=}1.3{\times}10^{-5}$), demonstrating that the causal signal is harvestable
under realistic constraints.
Against matched lightweight baselines, \bci-E1 leads the competitor matrix by 17+~pp,
outperforming Resample-Lite (60.6\%), CoVe-Lite (57.2\%), and Self-Correct (49.2\%).

\paragraph{Contributions.} We summarize our contributions as:
\begin{itemize}[leftmargin=1.5em,itemsep=1pt,topsep=2pt]
  \item We introduce \bci{}, the first systematic framework for \emph{premise-level}
    inference-time error correction in VLMs, explicitly distinguishing epistemic
    (premise) from logical (reasoning) failures.
  \item We provide causal evidence---via targeted \vs random ablation---that premise
    correction is a specific and effective intervention axis, with a $+35$~pp gap
    over random belief deletion at matched removal budget.
  \item We define an intervention policy family spanning from an oracle upper bound
    (E1: $+18.0$ pp, $+16.2$ pp at $n{=}2000$) to a practical confidence-gated
    variant (E12b: $+3.45$ pp at $n{=}2000$, $p{=}1.3{\times}10^{-5}$) suitable
    for deployment.
  \item We characterize verifier behavior across three profiles and four noise
    levels ($0\%$--$20\%$), showing that the pipeline degrades gracefully and that
    the 33.3\% premise-error rate is a conservative lower bound.
  \item We release a reproducible codebase with versioned verifier profiles,
    YAML-driven experiment configs, and full run manifests (git hash, code hash,
    config hash, random seed).
\end{itemize}
